problem:  
Let \((M^8,g)\) be a compact Riemannian spin manifold with holonomy contained in \(\mathrm{Spin}(7)\), and let \(D_g\) denote its Dirac operator. Since \(M\) admits a covariantly constant spinor, \(0\) is an eigenvalue of \(D_g\). Assume further that the next positive eigenvalue \(\lambda_1(g)\) of \(D_g^2\) is simple.  

Let \(\mathcal{M}\) be the local moduli space of torsion‑free \(\mathrm{Spin}(7)\)-structures near \(g\), modeled by harmonic 4‑forms in the Cayley representation. For tangent vectors \(X,Y\in T_g\mathcal{M}\), consider the metric perturbation \(g_{s,t}=g+sX+tY+O(s^2,t^2)\) (using the identification between deformations of the Cayley 4-form and of the metric). Define
\[
\mathcal{Q}(X,Y)=\left.\frac{\partial^2}{\partial s\,\partial t}\right|_{s=t=0}\lambda_1(g_{s,t}),
\]
which is well‑defined under the simplicity assumption by analytic perturbation theory.  

**Problem.**  
Does the bilinear form \(\mathcal{Q}\) on \(T_g\mathcal{M}\) detect the *holonomy type* of \(g\)? Specifically:
1. Is \(\mathcal{Q}\) positive‑definite if and only if the holonomy of \(g\) is *exactly* \(\mathrm{Spin}(7)\) (and not a proper subgroup such as \(\mathrm{SU}(4)\))?  
2. More broadly, does the infinitesimal spectral behavior of the Dirac operator across the moduli space of Spin(7) metrics encode the reduction pattern of the holonomy group?  

In other words, can one *hear the full exceptional holonomy* from the second‑order spectral response of \(D_g^2\) along torsion‑free deformations?  

Why is it a "good" problem:  
This problem formulates a precise spectral-geometric test for exceptional holonomy within the natural moduli space of torsion‑free Spin(7)-metrics. It improves conceptual clarity by working inside the well‑defined moduli space (so perturbations preserve the parallel spinor condition) and using analytic perturbation theory for a simple eigenvalue. The question lies at a deep intersection of spectral analysis, deformation theory of calibrated geometries, and the algebraic representation structure of \(\mathrm{Spin}(7)\). If answered, it could yield a spectral criterion distinguishing manifolds with full \(\mathrm{Spin}(7)\) holonomy from those with further reducible holonomy, thus approaching the open challenge “can one hear the holonomy of a manifold?” It is simple to state, mathematically rigorous in setup, and conceptually profound—demonstrating the unity of spin geometry, spectral theory, and exceptional holonomy.